\section{From Late Antiquity to the Medieval Bible}

Jerome's books did not replace the older Latin Scriptures by decree or in a single generation. They entered a manuscript culture in which books were copied separately, familiar readings could be restored from memory, and a new exemplar might itself combine several textual ancestries. By the sixth and seventh centuries Jerome's translations and the related New Testament revision had gained broad prestige, but Old Latin forms remained in circulation and continued to enter Vulgate copies. The medieval Vulgate was therefore both a successful received text and a history of continuing local correction (Houghton 2016, 55--68).

\subsection{Pandects and the evidence of Codex Amiatinus}

The single-volume Bible, or pandect, made the collection appear materially unified. Cassiodorus described biblical pandects at Vivarium in the sixth century, but no complete copy from his library survives. Around 700 the community of Wearmouth--Jarrow in Northumbria produced three giant Bibles. The survivor, Codex Amiatinus, is the earliest extant complete Vulgate Bible. Abbot Ceolfrid intended it as a gift for Rome and set out with it in 716; the codex eventually reached Italy and is now Florence, Biblioteca Medicea Laurenziana, Amiatino 1 (Houghton 2016, 69--72; Library of Congress, Codex Amiatinus record).

Amiatinus is exceptionally important without being an autograph or a transparent copy of Jerome's lost Bible. Its books were assembled by Northumbrian scholars from more than one source. Its Gospel text reflects an Italian tradition, while other parts have distinct histories. The codex establishes what one learned community could gather about three centuries after Jerome. It cannot by itself establish every form that Jerome or the anonymous New Testament reviser wrote.

\evidenceframe{Codex Amiatinus, made at Wearmouth--Jarrow around 700 and preserved as Florence, Biblioteca Medicea Laurenziana, Amiatino 1.}{The earliest surviving complete Vulgate Bible already presents the diverse component books as one monumental pandect and supplies an early, unusually valuable textual witness.}{It postdates Jerome by roughly three centuries, is not copied from one surviving authorial exemplar, and should not be made the sole norm for every book.}

The codex also exposes a useful distinction. A Bible may be ``complete'' as an artifact while its individual books remain textually composite. Contents, order, prologues, layout, and verbal readings can each have a separate history. Physical wholeness is evidence about reception, not proof of single authorship.

\subsection{Carolingian correction was plural}

The Carolingian reforms increased both the demand for correct service books and the capacity to produce them. Charlemagne's \emph{Admonitio generalis} of 789 called for experienced scribes to copy Gospel books, Psalters, and missals carefully. The best-known biblical programs associated with Alcuin of York and Theodulf of Orl\'eans answered that need differently.

Alcuin reported a royal command to emend the Old and New Testaments and completed a pandect around 800. His chief aims included grammar, punctuation, readability, and consistent presentation. The great productivity of the Tours scriptorium, rather than a surviving empire-wide act making one text compulsory, gave the Alcuinian form its reach. Even Tours Bibles were not identical: teams could work from several exemplars, and later books added or altered contents and prologues. The attribution of many prefatory texts to Jerome helped make the whole collection appear to be his work (Houghton 2016, 81--85; British Library, Add MS 10546).

Theodulf compared a wider set of Latin witnesses, including Spanish and Alcuinian forms, and marked selected variants in margins. His approach was closer to an annotated comparison than to Alcuin's regularized school and liturgical book. Marginal alternatives could later be copied into the text they were meant to document. Only a small number of Theodulfian Bibles survive; British Library Add MS 24142 is one of the principal witnesses. Neither program created a single Carolingian ``official Vulgate,'' and their different purposes resist being collapsed into one revision (Houghton 2016, 85--86; British Library, Add MS 24142).

\begin{fourstable}{Program}{Primary concern}{Mechanism of influence}{Historical limit}
Alcuin and Tours & Correct Latin, punctuation, presentation, and complete pandects & High-volume, carefully organized production at Tours & Not one invariant text and not demonstrably imposed as the sole imperial edition.\\
Theodulf and Orl\'eans/Fleury & Comparison of differing Latin witnesses and preservation of selected variants & A small, learned family of three-column Bibles with marginal signs & Variants could migrate into the text; relatively few copies survive.\\
Local monastic and cathedral copying & Worship, reading, teaching, and available exemplars & Repeated reproduction, correction, and mixture & ``Local'' does not mean careless, and a geographical label does not define every reading.\\
Later medieval book trade & Portable whole Bibles and shared study tools & Professional workshops, common order, divisions, and paratext & Standard format did not create a single critically controlled verbal text.\\
\end{fourstable}

\subsection{Glossed Bibles, Paris Bibles, and reference systems}

From the twelfth century, the biblical text increasingly traveled with structured exposition. Glossed books surrounded Scripture with patristic and medieval commentary; masters and students needed stable ways to locate and compare passages. The so-called Paris Bible, widespread from the thirteenth century, answered this material and educational setting with a relatively compact pandect, a common order of books, a recurring set of prologues, running titles, and chapter divisions close to those still used. The type spread because Paris was a center of schools and commercial book production, not because one act promulgated its text (Houghton 2016, 101--106; Light 2012).

This standardization was powerful but selective. A shared chapter number allowed a lecturer, concordance, or cross-reference to point reliably to a place even when two copies differed in wording. The reference architecture could therefore become more uniform while the text remained variable. Modern verse numbers belong to the age of print and should not be projected into Jerome's codices or the earliest Paris Bibles.

The thirteenth-century \emph{correctoria} make the continuing problem visible. Dominican and Franciscan scholars compiled variant readings from Latin manuscripts, Greek and Hebrew evidence, and patristic citations. Some alternatives stood in a Bible's margin; others circulated in separate reference works. Competing correctoria reveal both serious source criticism and disagreement about how much a received Latin text should be altered. Roger Bacon's criticism of particular Dominican corrections is evidence of the debate, not proof that one medieval school possessed Jerome's pristine text (Houghton 2016, 106--109).

\subsection{What ``the medieval Vulgate'' can mean}

The phrase may name at least three different objects: the Latin Bible as the West's common scriptural medium; a regional or institutional family of manuscripts; or the exact text of a particular codex. These levels overlap but are not interchangeable. The Vulgate could be the normal language of theology, liturgy, preaching, and schooling while individual copies still differed, retained Old Latin readings, or followed distinct Psalters.

That combination explains both the tradition's stability and its variation. Stability arose from worship, commentary, recognizable wording, and increasingly common book design. Variation arose from the ordinary mechanisms of manuscript transmission and from deliberate correction toward different authorities. Printing would greatly narrow the range that an ordinary reader encountered, but it would first reproduce one late-medieval Vulgate rather than recover the fourth-century text automatically.
